% =====================================================================
%  Chapter 1 — Introduction: The Probabilistic Method
%  File: chapters/01_introduction.tex
%  Restyled to match Chapter 2 (boxed environments + Insight callouts).
%  Labels needed by later chapters are preserved (notably thm:szele and
%  the method of conditional expectations).
% =====================================================================
\providecommand{\E}{\mathbb{E}}
\providecommand{\Prob}{\mathbb{P}}
\providecommand{\ind}{\mathbf{1}}

\chapter{The Probabilistic Method}
\label{chap:intro}

The probabilistic method is a technique for proving that an object exists by
showing that a \emph{random} object has the desired property with positive
probability. One never points to the object; one merely argues that the
alternative --- that \emph{no} good object exists --- would force a
probability to be zero when it cannot be. The disorienting feature, the first
time one meets it, is that a question with no randomness in its statement is
answered by introducing randomness that the answer then discards.

Two facts do almost all the work in this chapter. The first is \emph{linearity
of expectation}, $\E[\sum_i X_i]=\sum_i\E[X_i]$, valid with no independence
assumption; the second is the \emph{union bound},
$\Prob(\bigcup_i E_i)\le\sum_i\Prob(E_i)$. From these we extract existence
through one elementary principle: if $\E[X]=\mu$ then some outcome has
$X\ge\mu$ and some has $X\le\mu$, and if $\Prob(\bigcup_i E_i)<1$ then some
outcome avoids every $E_i$. The rest of the chapter is variations on this
theme, sharpened twice --- by \emph{alteration} (fix the defects of a random
object) and by the \emph{Lov\'asz Local Lemma} (escape the union bound when
dependencies are sparse).

% =====================================================================
\section{A first example: large bipartite subgraphs}
% =====================================================================

\begin{theorem}{Large bipartite subgraph}{maxcut}
Every graph with $m$ edges contains a bipartite subgraph with at least $m/2$
edges.
\end{theorem}

\begin{proof}
Colour each vertex of $G=(V,E)$ black or white, independently and uniformly at
random, and let $E'$ be the set of edges joining opposite colours. The graph
$(V,E')$ is bipartite by construction. Each edge lands in $E'$ exactly when its
two endpoints receive different colours, which happens with probability
$1/2$, so by linearity $\E|E'|=\tfrac12|E|$. Some colouring therefore achieves
$|E'|\ge\tfrac12|E|$, and that colouring exhibits the required subgraph.
\end{proof}

\begin{insight}
This is the probabilistic method in miniature, and worth dwelling on because
every later argument is a descendant. Note what the proof does \emph{not} do:
it never produces a good bipartition, only certifies one exists. The bound
$m/2$ is also a genuine guarantee, not just an asymptotic --- and it is
essentially tight, since the complete graph $K_n$ has $\binom n2$ edges but its
largest cut has only $\lfloor n^2/4\rfloor\approx m/2$ of them. Read
algorithmically, the same statement says random colouring is a
$\tfrac12$-approximation for \textsc{Max-Cut}; the entire field of randomized
approximation grows from exactly this observation.
\end{insight}

% =====================================================================
\section{From existence to algorithm}
% =====================================================================

A first-moment proof asserts existence but offers no recipe. Often one can be
recovered, by replacing a single random choice for the whole object with a
sequence of \emph{conditional} choices.

\begin{remark}{The method of conditional expectations}{condexp}
Instead of colouring all vertices of $G$ at once, colour them one at a time. At
vertex $v_k$, the colours of $v_1,\dots,v_{k-1}$ are already fixed; compute the
conditional expectation of $|E'|$ under each of the two choices for $v_k$
(with the remaining vertices still random) and commit to whichever is larger.
Because the average of the two conditional expectations equals the expectation
before the choice, the running conditional expectation never decreases. It
starts at $m/2$ and ends, once all vertices are coloured, at the actual cut
size --- which is therefore at least $m/2$. The existence proof has become a
deterministic linear-time algorithm.
\end{remark}

The same idea is a workhorse in theoretical computer science. Consider
\textsc{Max-Sat}: given a boolean formula in conjunctive normal form, satisfy
as many clauses as possible.

\begin{theorem}{Expected satisfiability}{maxsat}
A CNF formula with $m$ clauses, where clause $C_j$ has $k_j$ distinct literals,
admits a truth assignment satisfying at least
$\displaystyle\sum_{j=1}^m\bigl(1-2^{-k_j}\bigr)$ clauses.
\end{theorem}

\begin{proof}
Set each variable to true or false independently with probability $1/2$. Clause
$C_j$ fails only if all $k_j$ of its literals evaluate to false, an event of
probability $2^{-k_j}$ by independence, so the indicator $Z_j$ of ``$C_j$
satisfied'' has $\E[Z_j]=1-2^{-k_j}$. Linearity gives
$\E\sum_j Z_j=\sum_j(1-2^{-k_j})$, and some assignment meets this average.
\end{proof}

\begin{insight}
For clauses of width $k_j\ge1$ this guarantees at least $m/2$ satisfied
clauses, and the bound improves as clauses lengthen --- width-$3$ clauses
alone already force a $\tfrac78$ fraction. Derandomized by conditional
expectations, this is Johnson's algorithm. The right way to see conditional
expectations is as a \emph{potential function}: the conditional mean is a
quantity you can always hold steady or raise by a local decision, so a greedy
walk down the tree of choices lands at an outcome no worse than the global
average. That reframing --- existence proof as potential function ---
recurs throughout combinatorial optimization.
\end{insight}

% =====================================================================
\section{Counting Hamilton paths: Szele's theorem}
% =====================================================================

A \emph{tournament} is an orientation of $K_n$, and a \emph{Hamilton path} is a
directed path through all $n$ vertices. The following 1943 result is often
called the first true application of the probabilistic method.

\begin{theorem}{Szele 1943}{szele}
There is a tournament on $n$ vertices with at least $\dfrac{n!}{2^{\,n-1}}$
Hamilton paths.
\end{theorem}

\begin{proof}
Orient each edge of $K_n$ by an independent fair coin. Fix an ordering
$\pi=(v_1,\dots,v_n)$ of the vertices and let $X_\pi$ indicate that $\pi$ is a
directed path, i.e.\ that all $n-1$ consecutive edges point forward. These are
independent fair events, so $\Prob(X_\pi=1)=2^{-(n-1)}$. Summing over all $n!$
orderings, linearity gives $\E\sum_\pi X_\pi=n!\,2^{-(n-1)}$, and some
tournament attains at least the average.
\end{proof}

\begin{insight}
The proof is pure existence: it guarantees a tournament rich in Hamilton paths
without naming one. It is also a first taste of the first moment's habit of
being right in form but loose in scale --- the true maximum is exponentially
larger in its base, $n!/(2-o(1))^n$, a fact requiring much heavier machinery
(Alon's 1990 permanent bound). We return to the minimization side and these
sharper asymptotics in Chapter~2; here Szele serves simply as the historical
seed of everything that follows.
\end{insight}

% =====================================================================
\section{Lower bounds for Ramsey numbers}
% =====================================================================

\begin{definition}{Ramsey number}{ramseynum}
The \emph{Ramsey number} $R(k,\ell)$ is the least $n$ such that every red/blue
edge-colouring of $K_n$ contains a red $K_k$ or a blue $K_\ell$. Ramsey (1929)
proved $R(k,\ell)$ is finite; for instance $R(3,3)=6$.
\end{definition}

\begin{remark}{A historical aside}{history}
Paul Erd\H{o}s (1913--1996), with roughly $1{,}500$ papers and over $500$
collaborators, is regarded as the father of the probabilistic method; the
Ramsey lower bound below comes from his 1947 paper. Frank Ramsey (1903--1930)
did foundational work in logic, philosophy, and economics before dying at
$26$. Because so much of the subject traces to Erd\H{o}s's circle, Hungarian
names abound; a pronunciation guide:
\begin{center}\small
\begin{tabular}{lll}
\toprule
\textbf{Spelling} & \textbf{Sounds like} & \textbf{Examples}\\
\midrule
s  & sh & Erd\H{o}s, Simonovits\\
sz & s  & Szemer\'edi, Lov\'asz\\
\bottomrule
\end{tabular}
\end{center}
\end{remark}

\begin{theorem}{Erd\H{o}s 1947}{ramsey1}
If $\binom nk 2^{\,1-\binom k2}<1$, then $R(k,k)>n$.
\end{theorem}

\begin{proof}
Colour the edges of $K_n$ red/blue independently and uniformly. For each
$k$-set $S$ let $A_S$ be the event that $S$ is monochromatic; the
$\binom k2$ edges inside $S$ must agree, so $\Prob(A_S)=2^{\,1-\binom k2}$. By
the union bound,
\[
  \Prob\Bigl(\textstyle\bigcup_S A_S\Bigr)
  \le \sum_{S}\Prob(A_S)=\binom nk 2^{\,1-\binom k2}<1,
\]
so with positive probability no $A_S$ occurs --- a colouring with no
monochromatic $K_k$.
\end{proof}

\begin{remark}{Quantitative form}{ramsey1q}
Optimizing $n$ via Stirling gives
$R(k,k)>\bigl(\tfrac{1}{e\sqrt2}+o(1)\bigr)k\,2^{k/2}$. From the other side,
Erd\H{o}s--Szekeres (1935) showed $R(k,k)\le(4-o(1))^k$, and Campos,
Griffiths, Morris, and Sahasrabudhe (2023+) achieved the first exponential
improvement, $R(k,k)\le(4-c)^k$ for some $c>0$.
\end{remark}

\begin{insight}
The union bound is nothing but the first moment in disguise: $\sum_S\Prob(A_S)$
is the \emph{expected} number of monochromatic cliques, and ``expected number
$<1$'' forces ``some outcome has zero.'' The proof is also the standard
illustration of existence without construction --- a random colouring works
with overwhelming probability, yet no efficient way is known to certify that a
given colouring is good, a situation memorably described as finding hay in a
haystack. Pinning down an explicit construction matching the random bound is a
deep open problem with surprising ties to randomness extractors.
\end{insight}

\subsection*{Alteration}

\begin{theorem}{Ramsey lower bound via alteration}{ramsey_alt}
For all $k,n$, $\;R(k,k)>n-\binom nk 2^{\,1-\binom k2}$.
\end{theorem}

\begin{proof}
Colour $K_n$ at random as before, then \emph{delete} one vertex from each
monochromatic $K_k$; the result has no monochromatic $K_k$. If $X$ counts the
monochromatic cliques after step one, then $\E[X]=\binom nk 2^{\,1-\binom k2}$,
and we remove at most $X$ vertices, leaving at least $n-X$ with expectation
$n-\binom nk 2^{\,1-\binom k2}$. Some outcome meets this expectation.
\end{proof}

\begin{remark}{}{ramsey_altq}
Optimizing $n$ yields $R(k,k)>\bigl(\tfrac1e+o(1)\bigr)k\,2^{k/2}$, a factor
$\sqrt2$ better than the union-bound estimate.
\end{remark}

\begin{insight}
Alteration --- also called the deletion method --- is the first genuine upgrade
to the bare first moment: rather than demanding that a random object be perfect,
we accept a flawed one and pay to repair it. The accounting is the same first
moment, now applied to the \emph{cost of repair} (the expected number of
deletions) rather than to success outright. The payoff here is only a constant
factor, but the pattern --- build sloppily, then clean up --- is one of the
most broadly useful in the subject.
\end{insight}

\subsection*{The Lov\'asz Local Lemma}

To avoid bad events $E_1,\dots,E_m$ there are two crude regimes. The union
bound needs $\sum_i\Prob(E_i)<1$, which collapses once the events are
numerous. Full independence gives
$\Prob(\bigcap_i \overline{E_i})=\prod_i(1-\Prob(E_i))>0$, but genuine
independence is rare. The local lemma occupies the ground between: sparse
dependence suffices.

\begin{theorem}{Lov\'asz Local Lemma (variable model)}{lll}
Let $x_1,\dots,x_N$ be independent random variables and
$B_1,\dots,B_m\subseteq[N]$, with each event $E_i$ determined by
$\{x_j:j\in B_i\}$. If every $B_i$ meets at most $d$ other $B_j$ and
$\Prob(E_i)\le\dfrac{1}{(d+1)e}$ for all $i$, then with positive probability
none of the $E_i$ occur. (The constant $e$ is optimal; proof deferred to
Chapter~6.)
\end{theorem}

\begin{theorem}{Spencer 1977}{ramsey_lll}
If $\;e\bigl(\binom k2\binom{n}{k-2}+1\bigr)2^{\,1-\binom k2}\le 1$, then
$R(k,k)>n$.
\end{theorem}

\begin{proof}
Colour $K_n$ at random, one independent variable per edge, and let $E_S$ be the
event that the $k$-set $S$ is monochromatic, so $\Prob(E_S)=2^{\,1-\binom k2}$.
Two cliques share an edge precisely when their vertex sets overlap in at least
two vertices, and each $S$ has at most $\binom k2\binom{n}{k-2}$ such partners.
With $d=\binom k2\binom{n}{k-2}$, the hypothesis is exactly the local-lemma
condition $\Prob(E_S)\le 1/((d+1)e)$, so positively often no $E_S$ occurs.
\end{proof}

\begin{remark}{}{ramsey_lllq}
Optimizing gives $R(k,k)>\bigl(\tfrac{\sqrt2}{e}+o(1)\bigr)k\,2^{k/2}$, the best
lower bound known.
\end{remark}

\begin{insight}
The three Ramsey bounds form a clean ladder, each rung a factor $\sqrt2$ above
the last:
\[
  \tfrac{1}{e\sqrt2}\ \xrightarrow{\text{union bound}}\
  \tfrac{1}{e}\ \xrightarrow{\text{alteration}}\
  \tfrac{\sqrt2}{e}\ \ (k\,2^{k/2}).
\]
It is the same problem solved three times with progressively sharper tools, and
a faithful map of the method's hierarchy --- first moment, then repair, then
local independence. Tellingly, decades of effort have not moved the
\emph{leading constant} past $\sqrt2/e$ from below, even as the upper bound has
only recently cracked $4^k$ from above; the multiplicative gap in
$R(k,k)^{1/k}$ between $\sqrt2$ and $4-c$ remains one of the famous open
problems of the area.
\end{insight}

% =====================================================================
\section{Extremal set systems}
% =====================================================================

A set family $\mathcal F$ on ground set $[n]$ is an \emph{antichain} if no
member contains another. The extremal question --- how large can an antichain
be --- has a one-line probabilistic answer.

\begin{theorem}{Sperner 1928}{sperner}
An antichain of subsets of $[n]$ has at most $\binom{n}{\lfloor n/2\rfloor}$
members.
\end{theorem}

This follows from a stronger weighted statement.

\begin{theorem}{LYM inequality}{lym}
For an antichain $\mathcal F\subseteq 2^{[n]}$,
$\displaystyle\sum_{A\in\mathcal F}\binom{n}{|A|}^{-1}\le 1.$
\end{theorem}

\begin{proof}
Draw a uniform random permutation $\sigma$ of $[n]$ and form the maximal chain
$\emptyset\subset\{\sigma(1)\}\subset\{\sigma(1),\sigma(2)\}\subset\cdots\subset[n]$.
A fixed set $A$ appears in this chain exactly when $\sigma$ lists the elements
of $A$ first (in some order) and then the rest, which happens for
$|A|!\,(n-|A|)!$ of the $n!$ permutations; hence
$\Prob(A\in\text{chain})=\binom{n}{|A|}^{-1}$. The chain is itself a chain, so
it contains at most one member of the antichain $\mathcal F$ --- the events
``$A\in\text{chain}$'' for $A\in\mathcal F$ are disjoint, and their
probabilities sum to at most $1$.
\end{proof}

Sperner's bound drops out, since $\binom{n}{|A|}\le\binom{n}{\lfloor n/2\rfloor}$
makes each term at least $\binom{n}{\lfloor n/2\rfloor}^{-1}$, so
$|\mathcal F|\le\binom{n}{\lfloor n/2\rfloor}$.

\begin{insight}
The trick is a normalization: a uniformly random maximal chain meets each set
$A$ with probability $1/\binom{n}{|A|}$, so weighting each set by that
reciprocal turns the geometric fact ``a chain hits an antichain at most once''
into the analytic inequality $\sum 1/\binom{n}{|A|}\le1$. It is again a
first-moment argument, this time bounding the expected number of hits ($\le1$)
from above; equality forces $\mathcal F$ to be a single middle layer. The same
random-ordering device, suitably adapted, proves the next two theorems too ---
it is the unifying engine of probabilistic extremal set theory.
\end{insight}

\begin{theorem}{Bollob\'as two families theorem 1965}{bollobas}
Let $A_1,\dots,A_m$ and $B_1,\dots,B_m$ be finite sets with $A_i\cap B_i=\emptyset$
for all $i$ and $A_i\cap B_j\ne\emptyset$ for all $i\ne j$. Then
$\displaystyle\sum_{i=1}^m\binom{|A_i|+|B_i|}{|A_i|}^{-1}\le 1.$
\end{theorem}

\begin{proof}
Take a uniform random ordering of $\bigcup_i(A_i\cup B_i)$ and let $E_i$ be the
event that every element of $A_i$ precedes every element of $B_i$. Restricted
to $A_i\cup B_i$, this asks the $|A_i|$ elements of $A_i$ to come first among
$|A_i|+|B_i|$, so $\Prob(E_i)=\binom{|A_i|+|B_i|}{|A_i|}^{-1}$. The events are
pairwise disjoint: if $E_i$ and $E_j$ both held, the disjointness and crossing
hypotheses on $A_i,B_i,A_j,B_j$ would be contradicted. Disjoint events sum to
at most $1$.
\end{proof}

\begin{remark}{Why it is natural: transversals}{bollobastrans}
The uniform case $|A_i|=r,\,|B_i|=s$ reads $m\le\binom{r+s}{r}$, and it is
sharp. It also answers a question about transversals. Call $T$ a
\emph{transversal} of a family if it meets every member, write $\tau$ for the
minimum transversal size, and call a family $\tau$-critical if deleting any
member lowers $\tau$. The largest $\tau$-critical $r$-uniform family with
$\tau=s+1$ has exactly $\binom{r+s}{r}$ members: take the $A_i$ as the sets and
$B_i$ as an $s$-element transversal of the family with $A_i$ removed, and the
hypotheses of the theorem hold. Sperner and LYM are themselves special cases,
via $B_i=[n]\setminus A_i$.
\end{remark}

A family is \emph{intersecting} if any two of its members meet.

\begin{theorem}{Erd\H{o}s--Ko--Rado 1961}{ekr}
If $n\ge 2k$, an intersecting family of $k$-subsets of $[n]$ has at most
$\binom{n-1}{k-1}$ members.
\end{theorem}

\begin{proof}
Arrange $1,\dots,n$ uniformly at random around a circle and call a $k$-set
\emph{contiguous} if it forms an arc. A fixed $k$-set is contiguous with
probability $n/\binom nk$, so the expected number of contiguous members of
$\mathcal F$ is $n|\mathcal F|/\binom nk$. On the other hand, any single
circular arrangement contains at most $k$ contiguous members of an intersecting
family: fixing one contiguous member $A$, the only arcs meeting $A$ come in
disjoint pairs once $n\ge 2k$, so at most $k$ of them (including $A$) can
pairwise intersect. Hence $n|\mathcal F|/\binom nk\le k$, giving
$|\mathcal F|\le\tfrac kn\binom nk=\binom{n-1}{k-1}$.
\end{proof}

\begin{insight}
This is Katona's \emph{cycle method}: average over rotations rather than over
subsets, so a global counting problem becomes a tiny local one --- how many
arcs of an intersecting family fit on one circle. The hypothesis $n\ge2k$ is
not cosmetic; for $n<2k$ \emph{every} pair of $k$-sets intersects by
pigeonhole, the family $\binom{[n]}{k}$ is itself intersecting, and the bound
fails. The extremal family --- all $k$-sets through a fixed element --- shows
the answer $\binom{n-1}{k-1}$ is exactly right.
\end{insight}

% =====================================================================
\section{Two-colourable hypergraphs}
% =====================================================================

A $k$-uniform hypergraph is \emph{$2$-colourable} (has \emph{property B}) if its
vertices can be coloured with two colours leaving no edge monochromatic. Let
$m(k)$ be the fewest edges in a $k$-uniform hypergraph that is \emph{not}
$2$-colourable. The first values are $m(2)=3$, $m(3)=7$ (the Fano plane), and
$m(4)=23$ (found by computer search); $m(k)$ is unknown for every $k\ge5$. The
probabilistic method brackets it from both sides.

\begin{theorem}{Lower bound, Erd\H{o}s 1964}{hyper_lower}
$m(k)\ge 2^{k-1}$ for every $k\ge2$.
\end{theorem}

\begin{proof}
Suppose the hypergraph has $m<2^{k-1}$ edges and colour the vertices with a
uniform random $2$-colouring. A fixed $k$-edge is monochromatic with
probability $2\cdot 2^{-k}=2^{-(k-1)}$, so the expected number of monochromatic
edges is at most $2^{-(k-1)}m<1$. Some colouring therefore has none, i.e.\ the
hypergraph is $2$-colourable.
\end{proof}

\begin{theorem}{Upper bound, Erd\H{o}s 1964}{hyper_upper}
$m(k)=O(k^2 2^k)$.
\end{theorem}

\begin{proof}
Put $n=k^2$ vertices and build $H$ by choosing $m$ edges uniformly at random
(with replacement) from all $k$-subsets. Fix a $2$-colouring splitting the
vertices into parts of sizes $a$ and $b=n-a$. A random edge avoids being
monochromatic unless it lands entirely in one part, and a short estimate gives
$\Prob(\text{edge monochromatic})\ge c\,2^{-k}$ for an absolute $c>0$. The $m$
edges are independent, so
\[
  \Prob(\chi\text{ is proper})\le(1-c\,2^{-k})^m\le e^{-c\,2^{-k}m},
\]
and a union bound over all $2^n$ colourings gives
$\Prob(H\text{ is }2\text{-colourable})\le 2^n e^{-c\,2^{-k}m}<1$ once
$m=O(k^2 2^k)$ (recall $n=k^2$). Some such $H$ is then non-$2$-colourable.
\end{proof}

\begin{insight}
These two proofs run the probabilistic method in opposite directions. The lower
bound shows a \emph{good} object exists (a proper colouring) by making the
expected number of defects small. The upper bound shows a \emph{bad} object
exists (a non-colourable hypergraph) by making every fixed colouring likely to
fail and union-bounding over all colourings. Recognizing which direction a
problem wants --- existence of structure versus existence of a hard instance ---
is half the battle. The current record narrows the gap to
$m(k)\gtrsim 2^k\sqrt{k/\log k}$, still proved probabilistically.
\end{insight}

% =====================================================================
\section{The list chromatic number of \texorpdfstring{$K_{n,n}$}{Knn}}
% =====================================================================

In \emph{list colouring}, every vertex receives its own list of allowed colours;
$G$ is \emph{$k$-choosable} if a proper colouring exists for \emph{every}
assignment of $k$-element lists, and the \emph{list chromatic number}
$\mathrm{ch}(G)$ is the least such $k$. Always $\chi(G)\le\mathrm{ch}(G)$, and
the gap can be large: $K_{3,3}$ is $2$-colourable but not $2$-choosable.

\begin{theorem}{Upper bound}{list_upper}
If $n<2^{k-1}$ then $K_{n,n}$ is $k$-choosable; equivalently
$\mathrm{ch}(K_{n,n})\le\lfloor\log_2(2n)\rfloor+1$.
\end{theorem}

\begin{proof}
Mark each colour ``L'' or ``R'' by an independent fair coin. On a left vertex,
discard the colours marked R; on a right vertex, discard those marked L. A
fixed vertex loses its whole list of size $k$ with probability $2^{-k}$, so by
the union bound some vertex is emptied with probability at most $2n\,2^{-k}<1$.
With positive probability every list keeps a colour, and since left and right
vertices now draw from disjoint colour pools (L versus R), assigning any
surviving colour is proper.
\end{proof}

\begin{theorem}{Lower bound}{list_lower}
If some non-$2$-colourable $k$-uniform hypergraph has $n$ edges, then $K_{n,n}$
is not $k$-choosable.
\end{theorem}

\begin{proof}
Let $H$ be such a hypergraph with edges $e_1,\dots,e_n$ and treat its vertices
as colours. Give the $i$-th left vertex and the $i$-th right vertex the list
$e_i$. A proper list colouring would pick one colour (vertex of $H$) per list
(edge of $H$); the left choices and right choices each induce a colour class
that must hit every edge, forcing a proper $2$-colouring of $H$ --- which does
not exist.
\end{proof}

Combining the two with the hypergraph bound $m(k)=O(k^2 2^k)$ gives the
asymptotics.

\begin{corollary}{}{list_cor}
$\mathrm{ch}(K_{n,n})=(1+o(1))\log_2 n.$
\end{corollary}

\begin{insight}
The lower bound is the striking half: list-colourability of a graph is governed
by a $2$-colourability question about a hypergraph, so the difficulty of
property B transfers wholesale to choosability. Unlike the ordinary chromatic
number, $\mathrm{ch}$ keeps growing with average degree --- Saxton and Thomason
(2015) proved $\mathrm{ch}(G)>(1+o(1))\log_2 d$ for average degree $d$, with
$K_{n,n}$ showing this is sharp, via the method of \emph{hypergraph
containers}, one of the major modern developments in the area.
\end{insight}

% =====================================================================
\section*{Exercises}
\addcontentsline{toc}{section}{Exercises}
% =====================================================================

\begin{enumerate}[itemsep=4pt]
  \item Verify the asymptotic calculations behind the Ramsey lower bounds:
        \begin{enumerate}
          \item the largest $n$ with $\binom nk 2^{1-\binom k2}<1$ is
                $\bigl(\tfrac{1}{e\sqrt2}+o(1)\bigr)k\,2^{k/2}$;
          \item the maximum of $n-\binom nk 2^{1-\binom k2}$ over $n$ is
                $\bigl(\tfrac1e+o(1)\bigr)k\,2^{k/2}$;
          \item the largest $n$ with
                $e(\binom k2\binom{n}{k-2}+1)2^{1-\binom k2}<1$ is
                $\bigl(\tfrac{\sqrt2}{e}+o(1)\bigr)k\,2^{k/2}$.
        \end{enumerate}
  \item Prove that if some $p\in[0,1]$ satisfies
        $\binom nk p^{\binom k2}+\binom nt(1-p)^{\binom t2}<1$, then
        $R(k,t)>n$. Deduce $R(4,t)\ge c\,(t/\log t)^{3/2}$.
  \item Let $G$ have $n$ vertices and $m$ edges. Prove $K_n$ is a union of
        $O(n^2(\log n)/m)$ isomorphic copies of $G$ (not necessarily
        edge-disjoint).
  \item Prove there is an absolute $C>0$ such that every $n\times n$ matrix with
        distinct real entries can have its rows permuted so that no column
        contains an increasing subsequence of length $\ge C\sqrt n$.
  \item (Generalized Sperner.) If $\mathcal F\subseteq 2^{[n]}$ contains no
        chain $A_1\subsetneq\cdots\subsetneq A_{k+1}$, prove $|\mathcal F|$ is at
        most the union of the $k$ middle layers of the Boolean lattice.
  \item Let $G$ have $n\ge10$ vertices and suppose adding any edge creates a new
        $K_{10}$. Prove $G$ has at least $8n-36$ edges. \textit{(Hint: Bollob\'as.)}
  \item Let $k\ge4$ and let $H$ be a $k$-uniform hypergraph with at most
        $4^{k-1}/3^k$ edges. Prove its vertices can be $4$-coloured so that
        every edge sees all four colours.
  \item For $\mathcal F\subseteq 2^{[n]}$ and $A\subseteq[n]$ write
        $\mathcal F|_A=\{S\cap A:S\in\mathcal F\}$. Prove that for every $n,k$
        there is $\mathcal F$ with $|\mathcal F|=O(k2^k\log n)$ such that
        $\mathcal F|_A=2^A$ for every $k$-set $A$.
  \item Let $A_i$ be $r$-sets and $B_i$ be $s$-sets with $A_i\cap B_i=\emptyset$,
        and for each $i\ne j$ at least one of $A_i\cap B_j$, $A_j\cap B_i$
        nonempty. Prove $m\le(r+s)^{r+s}/(r^r s^s)$.
  \item $\bigstar$ Show that every non-$2$-colourable $n$-uniform hypergraph has
        at least $\tfrac n2\binom{2n-1}{n}$ unordered pairs of edges meeting in
        exactly one vertex.
\end{enumerate}
